\documentclass[11pt,a4paper]{article}
\usepackage[utf8]{inputenc}
\usepackage[T1]{fontenc}
\usepackage[english]{babel}
\usepackage{geometry}
\usepackage{microtype}
\usepackage{hyperref}
\usepackage{listings}
\usepackage{xcolor}
\usepackage{fancyhdr}

\geometry{margin=2.5cm, bottom=3cm}

% Configuration des liens hypertexte (aspect professionnel sans cadres colorés voyants)
\hypersetup{
    colorlinks=true,
    linkcolor=blue!50!black,
    citecolor=blue!50!black,
    urlcolor=blue!50!black
}

% Configuration des en-têtes et pieds de page professionnels
\pagestyle{fancy}
\fancyhf{}
\fancyfoot[L]{\footnotesize\ttfamily rostand.assonfo-demaboi@thga.stud.de}
\fancyfoot[C]{THGA Bochum}
\fancyfoot[R]{\thepage}
\renewcommand{\headrulewidth}{0pt}
\renewcommand{\footrulewidth}{0.4pt}

\lstset{
    basicstyle=\ttfamily\small\color{white},
    backgroundcolor=\color{black},
    rulecolor=\color{black},
    frame=single,
    breaklines=true,
    captionpos=b,
    xleftmargin=0pt,
    xrightmargin=0pt
}

\title{\textbf{User Manual -- Football Tournament Management System}}
\author{Rostand Cedrix Assonfo Demaboi}
\date{Summer Term 2026}

\begin{document}

\maketitle

\noindent\textbf{Technische Hochschule Georg Agricola, Bochum} \\
\textbf{Module:} Introduction to Database Management Systems \\
\textbf{Dozent:} Stephan Bökelmann (\href{mailto:sboekelmann@ep1.rub.de}{sboekelmann@ep1.rub.de})

\vspace{1cm}

\section{Introduction and System Overview}

The \emph{Football Tournament Management System} is an application designed to organize, track, and administer football competitions. The system relies on a decoupled architecture:
\begin{itemize}
    \item \textbf{Backend}: A REST API developed with FastAPI and connected to a relational database.
    \item \textbf{Frontend}: A desktop graphical interface developed in Python Tkinter, handling user interaction.
    \item \textbf{Orchestration}: Containerized deployment via Docker Compose ensuring infrastructure portability and stability.
\end{itemize}

Two main user roles are supported by the system:
\begin{itemize}
    \item \textbf{Organizer Mode}: Full access to write operations (creation, modification, and deletion of tournaments, teams, and matches) via an authentication key (\texttt{X-API-Key}).
    \item \textbf{Viewer Mode}: Read-only mode accessible to all participants, allowing them to view schedules, scores, and real-time updated rankings.
\end{itemize}

\section{Backend Deployment and Launch (Docker Compose)}

The backend server and database run via Docker Compose to ensure runtime environment integrity.

\subsection*{Deployment Procedure}
\begin{enumerate}
    \item Open a terminal and navigate to the project root directory:
\begin{lstlisting}
cd ~/football-tournament-manager
\end{lstlisting}
    \item Launch the containers in the background:
\begin{lstlisting}
docker compose up --build -d
\end{lstlisting}
    \item Check the status of active services:
\begin{lstlisting}
docker compose ps
\end{lstlisting}
\end{enumerate}

\section{Frontend Launch and Connection}

Once the backend server is operational on port 8000, the client graphical interface can be started.

\subsection*{Launch Commands}
\begin{lstlisting}
cd frontend
source venv/bin/activate
python3 -m football_tournament
\end{lstlisting}

\subsection*{Startup Connection Dialog}
At startup, the application presents a window requesting connection details:
\begin{itemize}
    \item \textbf{API URL}: Specify the server address (default: \url{http://localhost:8000}).
    \item \textbf{X-API-Key}:
    \begin{itemize}
        \item Enter the administrator key (\texttt{mykey}) to enable Organizer Mode.
        \item Leave the field blank to access Viewer Mode directly.
    \end{itemize}
\end{itemize}

\section{Interface Navigation and Usage}

The main interface is structured around a tabbed dashboard allowing seamless navigation between the system's various entities.

\subsection*{Tournaments Tab}
\begin{itemize}
    \item Displays the complete list of registered competitions.
    \item Allows viewing start dates, end dates, and the current status of each tournament.
\end{itemize}

\subsection*{Teams Tab}
\begin{itemize}
    \item Lists participating clubs with their unique identifier, official name, and coach/manager name.
\end{itemize}

\subsection*{Matches Tab}
\begin{itemize}
    \item \textbf{Global Tournament Selector}: Located at the top of the tab (active in both modes), it allows instantly filtering matches to show only those of a specific tournament or selecting "All Tournaments".
    \item \textbf{Score Format}: Planned matches with scores not yet defined explicitly display a dash (-). Validated scores are displayed in the \texttt{Home - Away} format.
    \item \textbf{Referee}: A dedicated column indicates the name of the referee assigned to the fixture.
\end{itemize}

\subsection*{Standings Tab}
\begin{itemize}
    \item Dynamically calculates team rankings per tournament based on aggregated match results (points, wins, draws, losses, goals for, goals against, and goal difference).
\end{itemize}

\section{Administrative Operations and Logout Guide}

Users authenticated as organizers have additional action buttons in the management tabs:
\begin{itemize}
    \item \textbf{Scheduling a Match}: Click the \texttt{+ Schedule Match} button to open the creation form, select the tournament, home and away teams, date, and save via \texttt{[SAVE MATCH]}.
    \item \textbf{Updating a Score}: Select a fixture in the matches table and click \texttt{Update Score} to modify the recorded goals.
    \item \textbf{Deletion}: Use the \texttt{Delete} button to remove an obsolete team or match from the database. All deletions strictly respect database referential integrity constraints.
    \item \textbf{Synchronization}: In case of external changes, use the \texttt{Refresh Data} button located in the upper toolbar to instantly refresh the local display from the API.
    \item \textbf{Logout}: To terminate an active session or switch user roles (e.g., from Organizer Mode back to Viewer Mode), simply close the application window or use the logout/disconnect option available in the connection panel. Reopening the app will prompt you again for credentials.
\end{itemize}

\end{document}